\section*{Capítulo \ref{ch:b2:aromatic-substitution} --- Aromaticidade e substituição eletrofílica aromática}

\begin{solution}{exo:b2:aromatic-substitution:1}
\ce{HNO3 + 2H2SO4 -> NO2+ + H3O+ + 2HSO4-}: o ácido nítrico é protonado e depois perde água.
\end{solution}

\begin{solution}{exo:b2:aromatic-substitution:2}
O tolueno dá 2-bromotolueno e 4-bromotolueno: a metila orienta para as
posições orto e para. O nitrobenzeno dá 3-bromonitrobenzeno, mais lentamente: o
grupo nitro orienta para meta.
\end{solution}

\begin{solution}{exo:b2:aromatic-substitution:3}
\ce{-CH3}: ativante, o/p. \ce{-OH}: fortemente ativante, o/p. \ce{-Cl}: desativante, o/p.
\ce{-NO2}: fortemente desativante, meta. \ce{-COCH3}: desativante, meta. \ce{-NHCOCH3}:
ativante (moderadamente), o/p.
\end{solution}

\begin{solution}{exo:b2:aromatic-substitution:4}
Acetofenona (1-feniletan-1-ona), \ce{C6H5COCH3}, por meio do \omterm{def:b2:aromatic-substitution:friedel-crafts}{íon acílio} \ce{CH3CO+}.
\end{solution}

\begin{solution}{exo:b2:aromatic-substitution:5}
O carbocátion primário (ou seu complexo) se rearranja por migração de hidreto no
cátion isopropila secundário, que alquila o anel. Propilbenzeno: \omterm{def:b2:aromatic-substitution:friedel-crafts}{acilação de
Friedel--Crafts} com cloreto de propanoíla (sem rearranjo), depois redução do
\ce{C=O} da cetona a \ce{CH2}.
\end{solution}

\begin{solution}{exo:b2:aromatic-substitution:6}
O grupo \ce{-OH} é um forte doador mesomérico: o anel fica tão ativado que o
bromo molecular é um eletrófilo suficiente, e as três posições orto e para são
substituídas.
\end{solution}

\begin{solution}{exo:b2:aromatic-substitution:7}
3-Bromonitrobenzeno: nitrar e depois bromar (o nitro orienta para meta).
4-Bromonitrobenzeno: bromar e depois nitrar (o bromo orienta para orto/para), e
separar o isômero para do orto.
\end{solution}

\begin{solution}{exo:b2:aromatic-substitution:8}
A mistura nitrante oxida a anilina, muito rica em elétrons, e protona seu
nitrogênio: \ce{-NH3+} é um \omterm{def:b2:aromatic-substitution:directing}{grupo desativante} e \omterm{def:b2:aromatic-substitution:directing}{orientador meta}. Na acetanilida,
o par não ligante do nitrogênio é compartilhado com a carbonila: menos básico e
menos ativante, ele não é protonado nem oxidado, e ainda orienta para
orto/para, sobretudo para, por razões estéricas.
\end{solution}

\begin{solution}{exo:b2:aromatic-substitution:9}
Sulfonar o tolueno (sobretudo em para, pois o grupo volumoso evita o orto);
bromar: o bromo entra em orto em relação à metila (a posição para está
bloqueada, e os dois grupos favorecem essa posição); remover o grupo \ce{SO3H}
aquecendo em ácido diluído.
\end{solution}

\begin{solution}{exo:b2:aromatic-substitution:10}
Os dois grupos orientam para orto/para; suas posições para estão ocupadas um
pelo outro. A metoxila, ativante mais forte, vence: nitração em orto em relação
ao \ce{-OCH3} (posição 2), dando o 4-metil-2-nitroanisol.
\end{solution}

\begin{solution}{exo:b2:aromatic-substitution:11}
Cada grupo nitro desativa fortemente o anel: a segunda nitração exige
condições mais fortes que a primeira, e a terceira (em um anel que leva dois
grupos nitro e apenas a metila, fracamente ativante), muito mais fortes ainda.
\end{solution}

\begin{solution}{exo:b2:aromatic-substitution:12}
Ácido 4-nitrobenzoico: nitrar o tolueno, separar o isômero para, oxidar a
metila a \ce{-COOH} (permanganato a quente). Ácido 3-nitrobenzoico: oxidar
primeiro o tolueno a ácido benzoico e depois nitrar (o \ce{-COOH} orienta para
meta).
\end{solution}

\begin{solution}{pb:b2:aromatic-substitution:1}
\textbf{1.} \ce{HNO3 + 2H2SO4 -> NO2+ + H3O+ + 2HSO4-}.
\textbf{2.} \ce{O=N+=O}: linear, isoeletrônico com o \ce{CO2}.
\textbf{3.} O nitrogênio se liga ao carbono do anel em para do \ce{CH3}, que se
torna tetraédrico (levando H e \ce{NO2}).
\textbf{4.} A carga positiva nos dois carbonos em orto do carbono atacado e no
carbono em para dele, que leva a metila.
\textbf{5.} A primeira, a formação do \omterm{def:b2:aromatic-substitution:seAr}{intermediário de Wheland}.
\textbf{6.} Uma base do meio: \ce{HSO4-} ou a água.
\textbf{7.} Duas orto, duas meta, uma para.
\textbf{8.} 2 : 2 : 1, isto é, 40~\% orto, 40~\% meta, 20~\% para.
\textbf{9.} Para o ataque em meta, a carga positiva nunca fica no carbono que
leva a metila; nenhuma estrutura é favorecida pela metila.
\textbf{10.} Uma estrutura de ressonância tem a carga no carbono que leva o
\ce{CH3}: um carbocátion terciário, estabilizado pelo efeito indutivo e pela
\omterm{def:b2:fragment-orbitals:hyperconjugation}{hiperconjugação} da metila.
\textbf{11.} Orto e para são favorecidas (96~\% contra 60~\% ao acaso); o meta
está quase ausente.
\textbf{12.} Por posição: orto $58/2 = 29~\%$, para 38~\%: para é a posição mais
reativa, pois as posições orto são ligeiramente impedidas pela metila.
\textbf{13.} Mais depressa: a metila ativa o anel.
\textbf{14.} O 4-nitrotolueno (funde a \qty{52}{\degreeCelsius}).
\textbf{15.} Resfriar a mistura bruta: o isômero para cristaliza e é filtrado; o
líquido retém os isômeros orto e meta.
\textbf{16.} Uma molécula simétrica se empacota melhor em um cristal: mais
interações por volume, um ponto de fusão mais alto.
\textbf{17.} Por \omterm{def:b2:liquid-vapour-diagrams:fractional-distillation}{destilação fracionada} sob pressão reduzida, ou por
cromatografia; ou por um resfriamento maior (o isômero meta funde a
\qty{16}{\degreeCelsius}).
\textbf{18.} Os isômeros têm pontos de ebulição próximos (mesma fórmula,
polaridade semelhante).
\textbf{19.} \ce{C7H8}: \qty{92.14}{g/mol}; \ce{C7H7NO2}: \qty{137.14}{g/mol}.
\textbf{20.} $100/92.14 = \qty{1.085}{mol}$ de tolueno; $0.90 \times 1.085 = \qty{0.977}{mol}$ de
nitrotoluenos.
\textbf{21.} Total $0.977 \times 137.14 = \qty{134.0}{g}$: orto \qty{77.7}{g}, meta
\qty{5.4}{g}, para \qty{50.9}{g}.
\textbf{22.} Uma segunda nitração (dinitrotoluenos) e a oxidação do grupo metila.
\textbf{23.} Oxidar a metila a \ce{-COOH} com permanganato alcalino a quente e
depois acidificar.
\textbf{24.} O ensaio dá $\boldsymbol{\approx \qty{51}{g}}$ de 4-nitrotolueno.
\end{solution}
